% Colour convention used in every schematic of the course
\documentclass[tikz,border=8pt]{standalone}
\usepackage{xfab-tikz}
\begin{document}
\begin{tikzpicture}[x=1cm,y=1cm]
\def\row#1#2#3{\node[lab, anchor=west] at (1.5,#1) {\textbf{#2}\\[-1pt]{\color{inksoft}\itshape #3}};}
% --- matter and tools
\fill[feed] (0,3.8) rectangle (0.5,4.4); \fill[ifeed] (0.6,3.8) rectangle (1.1,4.4);
                                                  \row{4.1}{initial material: kitchen, industry}{matière initiale · Ausgangsstoff}
\fill[part] (0,2.8) rectangle (0.5,3.4); \fill[ipart] (0.6,2.8) rectangle (1.1,3.4);
                                                  \row{3.1}{final material, the part: kitchen, industry}{pièce · Werkstück}
\fill[mouldx] (0,1.8) rectangle (1.1,2.4);         \row{2.1}{mould or die (carries the shape)}{moule, matrice · Form, Gesenk}
\fill[tool] (0,0.8) rectangle (1.1,1.4);           \row{1.1}{other tools}{outils · Werkzeuge}
\fill[fluid] (0,-0.2) rectangle (1.1,0.4);         \row{0.1}{carrier fluid (water, electrolyte, gas)}{fluide · Fluid}
% --- interactions
\begin{scope}[xshift=8.2cm]
\draw[force] (0,4.1) -- (1.1,4.1);                 \row{4.1}{force, motion}{effort, mouvement · Kraft, Bewegung}
\draw[heatin] (0,3.1) -- (1.1,3.1);                \row{3.1}{heat in or out}{chaleur · Wärme}
\draw[lighthalo] (0,2.1) -- (1.1,2.1); \draw[light] (0,2.1) -- (1.1,2.1);
                                                  \row{2.1}{radiation: laser, UV, electrons}{rayonnement · Strahlung}
\draw[elec] (0,1.1) -- (0.35,1.1) -- (0.5,1.3) -- (0.65,0.9) -- (0.8,1.1) -- (1.1,1.1);
                                                  \row{1.1}{electricity}{électricité · Elektrizität}
\fill[chem] (0.05,-0.05) circle (0.13); \fill[chem] (0.45,0.2) circle (0.13); \fill[chem] (0.9,0.0) circle (0.13);
                                                  \row{0.1}{chemistry: binder, reaction}{chimie · Chemie}
\end{scope}
\end{tikzpicture}
\end{document}
